\documentclass[12pt,a4paper]{article}
\usepackage[margin=2.5cm]{geometry}
\usepackage{amsmath,amssymb,graphicx,hyperref,natbib,booktabs,caption,xcolor}
\definecolor{ai4sgreen}{HTML}{0D3B2E}
\definecolor{ai4sgold}{HTML}{C8956C}

\title{\textbf{\textcolor{ai4sgold}{AI4S-Journal} Research Article}}
\author{Author Name\textsuperscript{1*}, Co-Author\textsuperscript{2} \
\textsuperscript{1}Department, Institution, City, Country \
\textsuperscript{2}Department, Institution, City, Country \
\textsuperscript{*}Corresponding author: email@institution.edu}
\date{}

\begin{document}
\maketitle

\begin{abstract}
Provide a concise abstract (200--300 words) summarizing the research question,
methodology, key findings, and significance. Must be self-contained.
\end{abstract}
\noindent\textbf{Keywords:} keyword1, keyword2, keyword3

\section{Introduction}
Introduce the research problem, review relevant literature, and state the
research question. End with a summary of contributions.

\section{Methods}
Describe your methodology in sufficient detail for reproducibility. Include
dataset description, model architecture, training protocol, hyperparameters,
and evaluation metrics with baselines.

\section{Results}
Present findings objectively. Include main experiments, comparison with
state-of-the-art, ablation studies, and statistical significance tests.

\section{Discussion}
Interpret results, discuss limitations, and contextualize within the field.

\section{Conclusion}
Summarize main findings and suggest future work.

\section*{Data and Code Availability}
A statement specifying where data and code supporting this study can be found.

\section*{Acknowledgments}
Acknowledge funding sources and contributors who do not meet authorship criteria.

\bibliographystyle{unsrt}
\begin{thebibliography}{99}
\bibitem{example1} Author, A. \& Author, B. (Year). Title. \textit{Journal}, \textit{Vol}(Issue), pages.
\end{thebibliography}
\end{document}
